% modul-5.tex - Modul 5: Fortgeschrittene KI-Integration

\startcomponent modul-5
\environment ../pyt-layout
\environment ../pyt-env

\chapter{Fortgeschrittene KI-Integration}

{\it Dauer: 1 Tag (4 Lektionen à 50 Minuten) • Voraussetzung: Module 1-4 abgeschlossen}

\startlernziele
Nach diesem Modul können Sie:
\startitemize[n,packed]
    \item LLM APIs (OpenAI, Anthropic) programmatisch nutzen
    \item RAG-Systeme (Retrieval-Augmented Generation) bauen
    \item AI Agents mit Tool-Nutzung entwickeln
    \item Embeddings und Vector Databases verwenden
    \item Ethische Aspekte verstehen und berücksichtigen
    \item Production-Ready AI-Apps erstellen
\stopitemize
\stoplernziele

% =============================================================================
% VORBEREITUNG
% =============================================================================

\section{Vorbereitung}

{\it Zeitaufwand: 2-3 Stunden vor dem Präsenzunterricht}

\subsection{API-Setup}

\startwarnbox
{\bf API-Keys sicher aufbewahren!}

API-Keys niemals in den Code schreiben oder committen. Verwenden Sie \code{.env}-Dateien:

\startbashcode
# .env (in .gitignore eintragen!)
OPENAI_API_KEY=sk-...
ANTHROPIC_API_KEY=sk-ant-...
\stopbashcode
\stopwarnbox

\subsubsection{OpenAI API einrichten}

\startitemize[n,packed]
    \item Account auf platform.openai.com erstellen
    \item API Key generieren
    \item In \code{.env}-Datei speichern
    \item Python-Package installieren: \code{pip install openai}
\stopitemize

\subsubsection{Kosten-Übersicht}

\starttabulate[|l|l|l|]
\HL
\NC {\bf Modell} \NC {\bf Input/1M Tokens} \NC {\bf Output/1M Tokens} \NC\NR
\HL
\NC GPT-4 Turbo \NC \$10 \NC \$30 \NC\NR
\NC GPT-3.5 Turbo \NC \$0.50 \NC \$1.50 \NC\NR
\NC Claude 3 Opus \NC \$15 \NC \$75 \NC\NR
\NC Claude 3 Sonnet \NC \$3 \NC \$15 \NC\NR
\HL
\stoptabulate

\starttippbox
{\bf Tipp für Entwicklung}

Für Entwicklung und Tests GPT-3.5 oder Claude Sonnet verwenden -- deutlich günstiger!
\stoptippbox

\subsection{Leseauftrag: LLM Basics}

\subsubsection{Was sind LLMs?}

{\bf Large Language Models} sind neuronale Netze mit Milliarden Parametern, trainiert auf riesigen Textmengen. Sie können:

\startitemize[packed]
    \item Text verstehen und generieren
    \item Code schreiben und erklären
    \item Fragen beantworten
    \item Texte übersetzen und zusammenfassen
\stopitemize

% =============================================================================
% PRAXIS
% =============================================================================

\section{Praxis}

{\it Präsenzunterricht: 4 Lektionen à 50 Minuten}

\subsection{Lektion 1: LLM APIs \& Prompt Engineering}

{\it Dauer: 50 Minuten}

\subsubsection{OpenAI API}

\startpythoncode
from openai import OpenAI
import os

# Client erstellen
client = OpenAI(api_key=os.getenv("OPENAI_API_KEY"))

# Chat Completion
response = client.chat.completions.create(
    model="gpt-3.5-turbo",
    messages=[
        {"role": "system", "content": "Du bist ein hilfreicher Assistent."},
        {"role": "user", "content": "Erklaere mir Python in 3 Saetzen."}
    ],
    temperature=0.7,
    max_tokens=500
)

print(response.choices[0].message.content)
\stoppythoncode

\subsubsection{Anthropic API (Claude)}

\startpythoncode
from anthropic import Anthropic
import os

client = Anthropic(api_key=os.getenv("ANTHROPIC_API_KEY"))

response = client.messages.create(
    model="claude-3-sonnet-20240229",
    max_tokens=1024,
    messages=[
        {"role": "user", "content": "Erklaere mir Python in 3 Saetzen."}
    ]
)

print(response.content[0].text)
\stoppythoncode

\subsubsection{Streaming Responses}

\startpythoncode
# Fuer lange Antworten - zeigt Token fuer Token
stream = client.chat.completions.create(
    model="gpt-3.5-turbo",
    messages=[{"role": "user", "content": "Schreibe eine Geschichte"}],
    stream=True
)

for chunk in stream:
    if chunk.choices[0].delta.content:
        print(chunk.choices[0].delta.content, end="")
\stoppythoncode

\subsection{Lektion 2: RAG-Systeme}

{\it Dauer: 50 Minuten}

\subsubsection{Was ist RAG?}

{\bf Retrieval-Augmented Generation} kombiniert:
\startitemize[n,packed]
    \item {\bf Retrieval:} Relevante Dokumente aus einer Datenbank holen
    \item {\bf Augmentation:} Dokumente als Kontext zum Prompt hinzufügen
    \item {\bf Generation:} LLM generiert Antwort basierend auf Kontext
\stopitemize

\subsubsection{RAG-Workflow}

\startpythoncode
# 1. Dokumente laden und in Chunks teilen
dokumente = lade_dokumente("./docs")
chunks = teile_in_chunks(dokumente, chunk_size=500)

# 2. Embeddings erstellen
embeddings = erstelle_embeddings(chunks)

# 3. In Vector-Datenbank speichern
vector_db.speichern(chunks, embeddings)

# 4. Bei Anfrage: Relevante Chunks suchen
frage = "Wie funktioniert die Login-Funktion?"
relevante_chunks = vector_db.suche(frage, top_k=3)

# 5. Kontext zum Prompt hinzufuegen
prompt = f"""Basierend auf folgenden Dokumenten:
{relevante_chunks}

Beantworte: {frage}"""

# 6. LLM antworten lassen
antwort = llm.generiere(prompt)
\stoppythoncode

\subsubsection{Embeddings}

\startpythoncode
# OpenAI Embeddings
from openai import OpenAI

client = OpenAI()

def get_embedding(text: str) -> list[float]:
    response = client.embeddings.create(
        model="text-embedding-3-small",
        input=text
    )
    return response.data[0].embedding

# Beispiel
embedding = get_embedding("Python ist eine Programmiersprache")
print(f"Dimension: {len(embedding)}")  # 1536
\stoppythoncode


\subsection{Lektion 3: AI Agents \& Tool Use}

{\it Dauer: 50 Minuten}

\subsubsection{Was sind AI Agents?}

AI Agents sind Systeme, die autonom Aufgaben lösen durch:
\startitemize[n,packed]
    \item {\bf Planung:} Aufgabe in Schritte zerlegen
    \item {\bf Tool-Nutzung:} Externe Funktionen aufrufen
    \item {\bf Iteration:} Ergebnisse prüfen und anpassen
    \item {\bf Selbst-Korrektur:} Fehler erkennen und beheben
\stopitemize

\subsubsection{Function Calling}

\startpythoncode
import json

# Tool-Definition
tools = [
    {
        "type": "function",
        "function": {
            "name": "get_weather",
            "description": "Holt das aktuelle Wetter",
            "parameters": {
                "type": "object",
                "properties": {
                    "city": {"type": "string", "description": "Stadt"}
                },
                "required": ["city"]
            }
        }
    }
]

# API-Call mit Tools
response = client.chat.completions.create(
    model="gpt-4-turbo",
    messages=[{"role": "user", "content": "Wie ist das Wetter in Zuerich?"}],
    tools=tools,
    tool_choice="auto"
)

# Tool-Call verarbeiten
if response.choices[0].message.tool_calls:
    tool_call = response.choices[0].message.tool_calls[0]
    args = json.loads(tool_call.function.arguments)

    # Funktion ausfuehren
    result = get_weather(args["city"])

    # Ergebnis zurueck an LLM
    ...
\stoppythoncode

\subsection{Lektion 4: Production \& Ethics}

{\it Dauer: 50 Minuten}

\subsubsection{Error Handling für APIs}

\startpythoncode
import time
from openai import RateLimitError, APIError

def sichere_api_anfrage(prompt: str, max_retries: int = 3) -> str:
    """API-Anfrage mit Retry-Logik."""
    for attempt in range(max_retries):
        try:
            response = client.chat.completions.create(
                model="gpt-3.5-turbo",
                messages=[{"role": "user", "content": prompt}]
            )
            return response.choices[0].message.content
        except RateLimitError:
            wait_time = 2 ** attempt  # Exponential Backoff
            print(f"Rate Limit - warte {wait_time}s")
            time.sleep(wait_time)
        except APIError as e:
            print(f"API Fehler: {e}")
            raise

    raise Exception("Max Retries erreicht")
\stoppythoncode

\subsubsection{Kosten-Kontrolle}

\startpythoncode
# Tokens zaehlen vor dem Senden
import tiktoken

def zaehle_tokens(text: str, model: str = "gpt-3.5-turbo") -> int:
    encoding = tiktoken.encoding_for_model(model)
    return len(encoding.encode(text))

# Kosten schaetzen
def schaetze_kosten(input_tokens: int, output_tokens: int) -> float:
    input_cost = (input_tokens / 1_000_000) * 0.50   # GPT-3.5
    output_cost = (output_tokens / 1_000_000) * 1.50
    return input_cost + output_cost
\stoppythoncode

\subsubsection{Ethik-Checkliste}

\startwarnbox
{\bf Ethische Aspekte beachten!}

\startitemize[packed]
    \item {\bf Bias:} Modelle können Vorurteile reproduzieren
    \item {\bf Datenschutz:} Keine persönlichen Daten an APIs senden
    \item {\bf Transparenz:} Nutzer informieren, dass KI verwendet wird
    \item {\bf Missbrauch:} Sicherheitsmassnahmen gegen böswillige Nutzung
    \item {\bf Umwelt:} Grosse Modelle haben hohen Energieverbrauch
\stopitemize
\stopwarnbox

% =============================================================================
% ÜBUNGEN
% =============================================================================

\section{Übungen}

\startaufgabenbox
{\bf Übung 1: Chatbot mit API} {\it (20 Min.)}

Einfachen Chatbot implementieren:
\startitemize[packed]
    \item Konversation mit Gedächtnis
    \item System-Prompt für Persönlichkeit
    \item Beenden mit "quit"
\stopitemize
\stopaufgabenbox

\blank[medium]

\startaufgabenbox
{\bf Übung 2: Dokument-Q\&A mit RAG} {\it (25 Min.)}

RAG-System für Textdatei:
\startitemize[packed]
    \item Text in Chunks teilen
    \item Embeddings erstellen
    \item Fragen beantworten lassen
\stopitemize
\stopaufgabenbox

\blank[medium]

\startaufgabenbox
{\bf Übung 3: Agent mit Tools} {\it (20 Min.)}

Agent mit Tool-Nutzung:
\startitemize[packed]
    \item Wetter-Tool (simuliert)
    \item Rechner-Tool
    \item Multi-Step Reasoning
\stopitemize
\stopaufgabenbox

\blank[medium]

\startaufgabenbox
{\bf Übung 4: Ethik-Analyse} {\it (15 Min.)}

Diskutieren Sie:
\startitemize[packed]
    \item Welche Risiken hat Ihre Anwendung?
    \item Wie können Sie Bias reduzieren?
    \item Wie schützen Sie Nutzerdaten?
\stopitemize
\stopaufgabenbox

% =============================================================================
% NACHBEARBEITUNG
% =============================================================================

\section{Nachbearbeitung}

{\it Hausaufgaben: 8-10 Stunden -- Abschlussprojekt}

\subsection{Aufgabe 1: RAG-System für Dokumentation}

Bauen Sie ein Q\&A-System:
\startitemize[packed]
    \item Eigene Dokumentation als Basis
    \item Vector-Datenbank (z.B. ChromaDB)
    \item CLI oder Web-Interface
    \item Quellenangabe in Antworten
\stopitemize

\subsection{Aufgabe 2: Multi-Agent-System}

Mehrere Agenten zusammenarbeiten lassen:
\startitemize[packed]
    \item Agent 1: Recherche
    \item Agent 2: Zusammenfassung
    \item Agent 3: Qualitätskontrolle
    \item Koordination zwischen Agenten
\stopitemize

\subsection{Aufgabe 3: Production-Ready AI-App}

Vollständige Anwendung erstellen:
\startitemize[packed]
    \item Error Handling
    \item Rate Limiting
    \item Kosten-Monitoring
    \item Logging
    \item Tests
\stopitemize

\subsection{Aufgabe 4: Abschlussprojekt}

Eigenes KI-Projekt fertigstellen:
\startitemize[packed]
    \item Vollständige Dokumentation
    \item README mit Installation
    \item Präsentation vorbereiten
    \item Code auf GitHub
\stopitemize

% =============================================================================
% MATERIALIEN
% =============================================================================

\section{Materialien}

\subsection{LLM API Cheat Sheet}

\starttabulate[|l|p(8cm)|]
\HL
\NC {\bf Aktion} \NC {\bf Code} \NC\NR
\HL
\NC Chat \NC \code{client.chat.completions.create(...)} \NC\NR
\NC Embedding \NC \code{client.embeddings.create(...)} \NC\NR
\NC Streaming \NC \code{stream=True} Parameter \NC\NR
\NC Tools \NC \code{tools=[...]} Parameter \NC\NR
\HL
\stoptabulate

\subsection{RAG Architecture}

\startitemize[n,packed]
    \item {\bf Load:} Dokumente laden
    \item {\bf Split:} In Chunks teilen (500-1000 Tokens)
    \item {\bf Embed:} Embeddings erstellen
    \item {\bf Store:} In Vector-DB speichern
    \item {\bf Retrieve:} Relevante Chunks suchen
    \item {\bf Generate:} Mit Kontext antworten
\stopitemize

\subsection{Projekt-Ideen}

\startitemize[packed]
    \item {\bf Dokumentations-Assistent} -- RAG über eigene Docs
    \item {\bf Code-Review-Agent} -- Pull Requests analysieren
    \item {\bf Research-Assistant} -- Informationen sammeln und zusammenfassen
    \item {\bf Personal Knowledge Base} -- Notizen mit semantischer Suche
\stopitemize

\blank[big]

\starttippbox
{\bf Erfolgskriterien für Modul 5 \& Kurs}

Sie haben den Kurs erfolgreich abgeschlossen, wenn Sie:
\startitemize[packed]
    \item LLM APIs programmatisch nutzen können
    \item Ein RAG-System implementiert haben
    \item AI Agents mit Tools bauen können
    \item Ethische Aspekte berücksichtigen
    \item Ein vollständiges Abschlussprojekt präsentieren können
\stopitemize
\stoptippbox

\blank[big]

\startinfobox
{\bf Herzlichen Glückwunsch!}

Sie haben alle 5 Module des Kurses "Python mit KI" abgeschlossen. Sie können nun:
\startitemize[packed]
    \item Python-Programme entwickeln
    \item KI-Tools effektiv nutzen
    \item Eigene KI-Anwendungen bauen
    \item Code-Qualität sicherstellen
    \item Mit APIs und Daten arbeiten
\stopitemize

Viel Erfolg bei Ihren zukünftigen Projekten!
\stopinfobox

\stopcomponent

